\documentclass[10pt,a4paper]{amsart}
\usepackage{amsmath,amssymb,amsthm,graphicx}
% Retained preamble for the Stacks Project chapter {CH} (shared project preamble).
% The chapter file pinned-original-source.tex only carries \input{{preamble}} on its first line,
% so no preamble text lives inside the chapter itself.  The Stacks Project uses one single shared
% preamble for every chapter: the snapshot stages it three times (sources/stacks-chow/preamble.tex,
% sources/stacks-etale-cohomology/preamble.tex and focused-seeds/AGP/sources/stacks-snapshot-20261011/
% preamble.tex) and those copies are identical in size.  The block below is copied mechanically,
% line for line, from sources/stacks-chow/preamble.tex lines 14-18 and 154-217: the Xy-pic setup and
% the native theorem declarations together with the mathematical macros.  Nothing was added, removed
% or redefined; the wrapper class and the amsmath/amssymb/amsthm/graphicx base packages are supplied
% by the assembler.
% For commutative diagrams we use Xy-pic
% Chapter: stacks-cohomology.tex -- Cohomology of Sheaves.
\usepackage[all]{xy}

% We use 2cell for 2-commutative diagrams.
\xyoption{2cell}

% Theorem environments.
%
\theoremstyle{plain}
\newtheorem{theorem}[subsection]{Theorem}
\newtheorem{proposition}[subsection]{Proposition}
\newtheorem{lemma}[subsection]{Lemma}

\theoremstyle{definition}
\newtheorem{definition}[subsection]{Definition}
\newtheorem{example}[subsection]{Example}
\newtheorem{exercise}[subsection]{Exercise}
\newtheorem{situation}[subsection]{Situation}

\theoremstyle{remark}
\newtheorem{remark}[subsection]{Remark}
\newtheorem{remarks}[subsection]{Remarks}

\numberwithin{equation}{subsection}

% Macros
%
\def\lim{\mathop{\mathrm{lim}}\nolimits}
\def\colim{\mathop{\mathrm{colim}}\nolimits}
\def\Spec{\mathop{\mathrm{Spec}}}
\def\Hom{\mathop{\mathrm{Hom}}\nolimits}
\def\Ext{\mathop{\mathrm{Ext}}\nolimits}
\def\SheafHom{\mathop{\mathcal{H}\!\mathit{om}}\nolimits}
\def\SheafExt{\mathop{\mathcal{E}\!\mathit{xt}}\nolimits}
\def\Sch{\mathit{Sch}}
\def\Mor{\mathop{\mathrm{Mor}}\nolimits}
\def\Ob{\mathop{\mathrm{Ob}}\nolimits}
\def\Sh{\mathop{\mathit{Sh}}\nolimits}
\def\NL{\mathop{N\!L}\nolimits}
\def\CH{\mathop{\mathrm{CH}}\nolimits}
\def\proetale{{pro\text{-}\acute{e}tale}}
\def\etale{{\acute{e}tale}}
\def\QCoh{\mathit{QCoh}}
\def\Ker{\mathop{\mathrm{Ker}}}
\def\Im{\mathop{\mathrm{Im}}}
\def\Coker{\mathop{\mathrm{Coker}}}
\def\Coim{\mathop{\mathrm{Coim}}}

% Boxtimes
%
\DeclareMathSymbol{\boxtimes}{\mathbin}{AMSa}{"02}

%
% Macros for moduli stacks/spaces
%
\def\QCohstack{\mathcal{QC}\!\mathit{oh}}
\def\Cohstack{\mathcal{C}\!\mathit{oh}}
\def\Spacesstack{\mathcal{S}\!\mathit{paces}}
\def\Quotfunctor{\mathrm{Quot}}
\def\Hilbfunctor{\mathrm{Hilb}}
\def\Curvesstack{\mathcal{C}\!\mathit{urves}}
\def\Polarizedstack{\mathcal{P}\!\mathit{olarized}}
\def\Complexesstack{\mathcal{C}\!\mathit{omplexes}}
% \Pic is the operator that assigns to X its picard group, usage \Pic(X)
% \Picardstack_{X/B} denotes the Picard stack of X over B
% \Picardfunctor_{X/B} denotes the Picard functor of X over B
\def\Pic{\mathop{\mathrm{Pic}}\nolimits}
\def\Picardstack{\mathcal{P}\!\mathit{ic}}
\def\Picardfunctor{\mathrm{Pic}}
\def\Deformationcategory{\mathcal{D}\!\mathit{ef}}
\usepackage{cleveref}
\crefname{originallemma}{original lemma}{original lemmas}
\Crefname{originallemma}{Original lemma}{Original lemmas}
\makeatletter
\newlabel{modules-lemma-abelian}{{01AG}{source}{}{sourcepointer-1}{}}
\newlabel{modules-lemma-abelian@cref}{{[originallemma][0][]01AG}{[1][0][]source}{}{}{}}
\makeatother
\crefname{originallemma}{original lemma}{original lemmas}
\Crefname{originallemma}{Original lemma}{Original lemmas}
\makeatletter
\newlabel{simplicial-lemma-homotopy-equivalence-s-Q}{{0FKK}{source}{}{sourcepointer-2}{}}
\newlabel{simplicial-lemma-homotopy-equivalence-s-Q@cref}{{[originallemma][0][]0FKK}{[1][0][]source}{}{}{}}
\makeatother
\crefname{originalsection}{original section}{original sections}
\Crefname{originalsection}{Original section}{Original sections}
\makeatletter
\newlabel{simplicial-section-dold-kan-cosimplicial}{{019H}{source}{}{sourcepointer-3}{}}
\newlabel{simplicial-section-dold-kan-cosimplicial@cref}{{[originalsection][0][]019H}{[1][0][]source}{}{}{}}
\makeatother
\title{Godement resolution of a sheaf of modules by flasque modules with stalks a homotopy equivalence}
\author{The Stacks Project Authors}
\date{stacks-project:cohomology.tex:local-snapshot-20261011; GNU FDL 1.2 or later}
\begin{document}
\maketitle
\noindent\textit{Editorial note.} All delivered mathematical text is the author's own, copied from the pinned Stacks Project chapter file cohomology.tex; no statement or proof step was rewritten, reordered or removed.\par\medskip
\noindent\textit{Editorial note.} The proof refers to a result 'mentioned in the introduction to this section'; the sentence is reproduced as written and the result in question is a navigation pointer carried as a stub, while the reduction that makes it applicable is inside the counted span.\par\medskip
\noindent\textit{Editorial note.} The counted spans carry 3 navigation stubs and no bibliography citation.\par\medskip
\noindent\textit{Editorial note.} The preamble is the chapter's retained minimal preamble (sources/stacks-cohomology/retained-preamble.tex lines 1-81), a verbatim excerpt of the native preamble.\par\medskip
\noindent\textit{Editorial note.} The label lemma-godement-resolution is defined only in cohomology.tex, so no cross-chapter restatement had to be weighed for this card.\par\medskip
\noindent\textit{Allowed standard prerequisites (editorial).}\begin{itemize}
\item Complex associated to a cosimplicial object (Dold-Kan, cosimplicial form). Simplicial, Section simplicial-section-dold-kan-cosimplicial, invoked at proof lines 7045-7046. Supplies the identification of the displayed complex with the one attached to the cosimplicial object f\_*f\textasciicircum{}* applied iteratively to F.
\item f\textasciicircum{}* is exact and exactness can be tested after pulling back. Modules, Lemma modules-lemma-abelian, invoked at proof lines 7054-7060. Supplies that a complex of O\_X-modules is exact if and only if its pullback is, which is what makes the stalkwise argument work.
\item Homotopy equivalence of a constant cosimplicial object with its associated cosimplicial object. Simplicial, results on the standard simplicial homotopy named in the introduction to the section and used at proof lines 7061-7070. Supplies the stalkwise homotopy equivalence statement of the lemma.
\end{itemize}\medskip
\ifdefined\hypertarget\hypertarget{sourcepointer-1}{}\fi
\noindent\textit{Source reference.} Original lemma 01AG: sources/stacks-chow/tags.txt, source lines 1689--1689. Pointer outside the counted span, invoked once for context; the load-bearing argument of the counted statement is the author's proof reproduced in full in the proof span.\par\medskip
\ifdefined\hypertarget\hypertarget{sourcepointer-2}{}\fi
\noindent\textit{Source reference.} Original lemma 0FKK: sources/stacks-chow/tags.txt, source lines 19465--19465. Pointer outside the counted span, invoked once for context; the load-bearing argument of the counted statement is the author's proof reproduced in full in the proof span.\par\medskip
\ifdefined\hypertarget\hypertarget{sourcepointer-3}{}\fi
\noindent\textit{Source reference.} Original section 019H: sources/stacks-chow/tags.txt, source lines 1652--1652. Pointer outside the counted span, invoked once for context; the load-bearing argument of the counted statement is the author's proof reproduced in full in the proof span.\par\medskip
\begin{lemma}
\label{lemma-godement-resolution}
Let $(X, \mathcal{O}_X)$ be a ringed space. For every sheaf of
$\mathcal{O}_X$-modules $\mathcal{F}$ there is a resolution
$$
0 \to
\mathcal{F} \to
f_*f^*\mathcal{F} \to
f_*f^*f_*f^*\mathcal{F} \to
f_*f^*f_*f^*f_*f^*\mathcal{F} \to \ldots
$$
functorial in $\mathcal{F}$ such that each term
$f_*f^* \ldots f_*f^*\mathcal{F}$ is a flasque
$\mathcal{O}_X$-module and such that for all $x \in X$ the
map
$$
\mathcal{F}_x[0] \to \Big(
(f_*f^*\mathcal{F})_x \to
(f_*f^*f_*f^*\mathcal{F})_x \to
(f_*f^*f_*f^*f_*f^*\mathcal{F})_x \to \ldots
\Big)
$$
is a homotopy equivalence in the category of complexes
of $\mathcal{O}_{X, x}$-modules.
\end{lemma}

\begin{proof}
The complex $f_*f^*\mathcal{F} \to  f_*f^*f_*f^*\mathcal{F} \to
f_*f^*f_*f^*f_*f^*\mathcal{F} \to \ldots$ is the complex associated
to the cosimplicial object with terms
$f_*f^*\mathcal{F}, f_*f^*f_*f^*\mathcal{F},
f_*f^*f_*f^*f_*f^*\mathcal{F}, \ldots$ described above, see
Simplicial, Section \ref{simplicial-section-dold-kan-cosimplicial}.
The augmentation gives rise to the map $\mathcal{F} \to f_*f^*\mathcal{F}$
as indicated. For any abelian sheaf $\mathcal{H}$ on $X_{disc}$ the
pushforward $f_*\mathcal{H}$ is flasque because $X_{disc}$
is a discrete space and the pushforward of a flasque sheaf is flasque.
Hence the terms of the complex are flasque $\mathcal{O}_X$-modules.

\medskip\noindent
If $x \in X_{disc} = X$ is a point, then $(f^*\mathcal{G})_x = \mathcal{G}_x$
for any $\mathcal{O}_X$-module $\mathcal{G}$. Hence $f^*$ is an exact functor
and a complex of $\mathcal{O}_X$-modules
$\mathcal{G}_1 \to \mathcal{G}_2 \to \mathcal{G}_3$
is exact if and only if
$f^*\mathcal{G}_1 \to f^*\mathcal{G}_2 \to f^*\mathcal{G}_3$
is exact (see Modules, Lemma \ref{modules-lemma-abelian}).
The result mentioned in the introduction to this section
proves the pullback by $f^*$ gives a homotopy equivalence from
the constant cosimplicial object $f^*\mathcal{F}$ to the
cosimplicial object with terms
$f_*f^*\mathcal{F}, f_*f^*f_*f^*\mathcal{F},
f_*f^*f_*f^*f_*f^*\mathcal{F}, \ldots$.
By Simplicial, Lemma \ref{simplicial-lemma-homotopy-equivalence-s-Q}
we obtain that
$$
f^*\mathcal{F}[0] \to \Big(
f^*f_*f^*\mathcal{F} \to
f^*f_*f^*f_*f^*\mathcal{F} \to
f^*f_*f^*f_*f^*f_*f^*\mathcal{F} \to \ldots
\Big)
$$
is a homotopy equivalence. This immediately implies the two remaining
statements of the lemma.
\end{proof}

\end{document}
